\questao{114}{corD}{Energia e calor}{Energia por distância: de kWh para joule}{Com conta}
\begin{enunciado}
% gerado por ferramentas/md2tex_enunciado.py a partir de questoes/q114.md (não editar à mão)
Uma fabricante de automóveis anunciou o plano de eletrificar toda a sua frota até o ano 2030. Determinado protótipo é parte fundamental dessa estratégia por ter autonomia de 1\,000\,km e um gasto energético de 1\,kWh para cada 10\,km, quando ocupado por quatro pessoas. O gráfico apresenta, para vários meios de transporte, o gasto energético por passageiro, considerado como um dos fatores determinantes para o desenvolvimento de políticas públicas de regulação dos meios de transporte.\par
\figq[0.74\linewidth]{q114_fig1.png}
\fonte{Disponível em: https://quatrorodas.abril.com.br. Acesso em: 2 dez. 2021 (adaptado).}
De acordo com o gráfico, quantos meios de transporte apresentam maior eficiência energética que esse protótipo?\par
\begin{altlinha}
\alth{A}{1}%
\alth{B}{2}%
\alth{C}{4}%
\alth{D}{5}%
\alth{E}{7}
\end{altlinha}

\end{enunciado}

\quer{quantos meios de transporte do gráfico gastam \textbf{menos} energia que o protótipo (o carro elétrico em teste), por metro e por passageiro.}

\begin{ideia}
Só dá para comparar coisas que estão \textbf{na mesma unidade}. É como comparar preços: um em reais e outro em dólares não se
comparam direto.

O gráfico usa \textbf{joule por metro por passageiro}: quantos joules são gastos para levar \textbf{uma pessoa} por
\textbf{um metro}. O protótipo foi dado de outro jeito: 1\un{kWh} para cada 10\un{km}, com 4 pessoas.
O trabalho da questão é passar o protótipo para a unidade do gráfico.

E atenção ao significado: \textbf{mais eficiente é quem gasta menos}. O próprio gráfico avisa: ``mais eficiente'' fica do
lado do zero.
\end{ideia}

\passoapasso
\npasso{Transforme 1\un{kWh} em joules.}
O quilowatt-hora é uma unidade de energia: é a energia de um aparelho de 1\un{kW} ligado durante 1 hora.
\begin{itemize}
  \item $1\un{kW} = 1\,000\un{W}$, e 1\un{W} é 1 joule por segundo. Então 1\un{kW} são 1\,000 joules por segundo.
  \item $1\un{h} = 60\un{min} = 60 \times 60 = 3\,600\un{s}$.
\end{itemize}
\begin{conta}
$1\un{kWh} = 1\,000 \times 3\,600 = 3\,600\,000\un{J}$
\end{conta}

\npasso{Transforme 10\un{km} em metros.}
\begin{conta}
$10\un{km} = 10 \times 1\,000 = 10\,000\un{m}$
\end{conta}

\npasso{Calcule o gasto do carro por metro.}
O carro gasta 3\,600\,000\un{J} para andar 10\,000\un{m}. Divida para saber o gasto em cada metro:
\begin{conta}
$\dfrac{3\,600\,000}{10\,000} = 360\un{J}$ por metro
\end{conta}
(Para dividir por 10\,000, corte quatro zeros: 3\,600\,000 vira 360.)

\npasso{Divida entre os passageiros.}
Esses 360\un{J} levam 4 pessoas. Para cada uma:
\begin{conta}
$\dfrac{360}{4} = 90\un{J}$ por metro por passageiro
\end{conta}

\npasso{Compare com o gráfico.}
Procure quem gasta \textbf{menos de 90}:
\begin{itemize}
  \item Velomóvel: 18 \ (menos que 90) \ \faCheck
  \item Bicicleta: 60 \ (menos que 90) \ \faCheck
  \item Pedonalidade (andar a pé): 132 \ (mais que 90) \ \faTimes
\end{itemize}
Daí para baixo no gráfico, todos gastam ainda mais. Então são \textbf{2} meios de transporte mais eficientes que o protótipo.

\resposta{B}{O protótipo gasta 90\un{J} por metro por passageiro. Só o velomóvel (18) e a bicicleta (60) gastam menos.}

\begin{outras}
\alt{C} 4 aparece para quem para em 360\un{J} por metro e esquece de dividir pelos 4 passageiros (18, 60, 132 e 200 ficam abaixo de 360).\par
\altl{A}, \altl{D} e \alt{E} não batem com a contagem: só duas barras (18 e 60) ficam abaixo de 90.
\end{outras}

\begin{pegadinha}[Onde muita gente erra]
\begin{itemize}
  \item \textbf{Esquecer de dividir pelos 4 passageiros.} O gráfico é \textbf{por passageiro}.
  \item \textbf{Usar a autonomia de 1\,000\un{km}} (a distância que o carro roda com uma carga). Esse dado não entra na conta.
  \item \textbf{Achar que ``maior eficiência'' é número maior.} É o contrário: mais eficiente é quem gasta \textbf{menos} energia.
\end{itemize}
\end{pegadinha}

\begin{guardar}
$1\un{kWh} = 3\,600\,000\un{J} = 3{,}6 \times 10^{6}\un{J}$. Essa conversão aparece muito no ENEM. Se esquecer, refaça:
1\,000 joules por segundo, durante 3\,600 segundos.
\end{guardar}
